\documentclass[11pt]{article}
\usepackage[margin=1.05in]{geometry}
\usepackage{amsmath,amssymb,amsthm,mathtools}
\usepackage{booktabs,enumitem,microtype}
\usepackage[dvipsnames]{xcolor}
\usepackage[colorlinks=true,linkcolor=MidnightBlue,citecolor=MidnightBlue,urlcolor=BrickRed]{hyperref}
\setlist{nosep,leftmargin=*}
\allowdisplaybreaks

\newtheorem{theorem}{Theorem}[section]
\newtheorem{proposition}[theorem]{Proposition}
\newtheorem{lemma}[theorem]{Lemma}
\newtheorem{corollary}[theorem]{Corollary}
\theoremstyle{remark}
\newtheorem{remark}[theorem]{Remark}

\newcommand{\T}{\mathbb T}
\newcommand{\N}{\mathbb N}
\newcommand{\Z}{\mathbb Z}
\newcommand{\e}{\mathrm e}
\newcommand{\E}{\mathcal E}
\newcommand{\M}{\mathfrak M}
\newcommand{\m}{\mathfrak m}
\newcommand{\Sg}{\mathfrak S}
\newcommand{\eps}{\varepsilon}
\newcommand{\1}{\mathbf 1}
\newcommand{\dd}{\,\mathrm d}
\newcommand{\norm}[1]{\left\lVert#1\right\rVert}
\newcommand{\abs}[1]{\left|#1\right|}
\DeclareMathOperator{\dist}{dist}
\DeclareMathOperator{\RePart}{Re}

\title{Prime Pairs at the Prime Number Theorem Threshold\\
\large Hardy--Littlewood for Almost Every Shift in Every
Interval of Length $X^{2/15+\varepsilon}$}
\author{Conor Grogan}
\date{20 September 2026}

\begin{document}
\maketitle

\begin{abstract}
Let
\[
 S_X(\alpha)=\sum_{X<n\le 2X}\Lambda(n)\e^{2\pi i n\alpha}.
\]
For every \(A,\eps>0\), we choose narrow neighborhoods of rational points with
logarithmically bounded denominators and prove
\[
 \sup_{|J|=1/H}\int_{J\setminus\mathfrak M}|S_X(\alpha)|^2\,d\alpha
 \ll_{A,\eps}X(\log X)^{-A}
\]
whenever \(H\ge X^{2/15+\eps}\).  It follows that, in every translated shift
interval of the stated length, the Hardy--Littlewood prime-pair formula
holds in mean square and for all but \(O(H(\log X)^{-A})\) individual
shifts, and the second moment of \(R_X\) matches its singular-series
prediction.  The same estimate gives the weighted Goldbach formula for
almost every target in \([X,X+H]\), together with a balanced version whose
summands lie in \((X,2X]\).  This replaces the exponent \(8/33\) in
Matom\"aki--Radziwi\l\l--Tao by \(2/15\).
\end{abstract}

\section*{Introduction}
The prime number theorem in short intervals counts one copy of the primes:
\[
 \psi(x+y)-\psi(x)\sim y.
\]
The Hardy--Littlewood problem correlates the primes with a translate:
\[
 R_X(h)=\sum_{X<n\le2X}\Lambda(n)\Lambda(n+h)
 \sim X\mathfrak S(h).
\]
The singular series \(\mathfrak S(h)\) is the product of the local congruence
factors.  At \(h=2\), this correlation contains the twin-prime problem.
The theorems below are averages over \(h\) and leave every fixed shift open.
Averaging in \(h\) lets us reach the \(2/15\) short-interval scale.

The current one-point theory gives the prime number theorem for almost all
intervals of length \(x^{2/15+\eps}\) \cite{GafniTao2026}.  Matom\"aki,
Radziwi\l\l, and Tao \cite{MRT2019} proved the Hardy--Littlewood prime-pair
asymptotic for almost all shifts in families of length
\(X^{8/33+\eps}\).  Theorem~\ref{thm:map} lowers the two-point threshold to
\(X^{2/15+\eps}\), so the one-point and two-point thresholds now agree.
For \(X=10^{30}\), the old and new shift scales are approximately
\(10^{7.27}\) and \(10^4\): about nineteen million shifts versus ten
thousand.

The estimate is uniform in the center of the shift interval: in every
translated interval of length \(X^{2/15+\eps}\), almost all shifts obey the
Hardy--Littlewood prediction.

\section{Statements}
Write $L=\log X$, $e(t)=\e^{2\pi i t}$, and identify
$\T=\mathbb R/\mathbb Z$.  For positive $P,Q$ set
\[
 \M(P,Q)=\bigcup_{q\le Q}\ \bigcup_{\substack{a\bmod q\\(a,q)=1}}
 \left\{\alpha\in\T:\left\|\alpha-\frac aq\right\|_{\T}\le\frac PX\right\},
 \qquad \m(P,Q)=\T\setminus\M(P,Q).
\]

\begin{theorem}[Uniform local minor-arc estimate]\label{thm:map}
For every $A>0$ and $\eps>0$ there are positive integers
$B=B(A,\eps)$, $D=D(A,\eps)$ and positive constants $C=C(A,\eps)$ and
$X_0=X_0(A,\eps)$ such that, for $X\ge X_0$, with
\[
 Q=L^B,\qquad P=L^D,
\]
one has, simultaneously for every real $H\ge X^{2/15+\eps}$,
\[
 \sup_{\substack{J\subset\T\ \mathrm{an\ arc}\\ |J|=1/H}}
 \int_{J\setminus\M(P,Q)}\abs{S_X(\alpha)}^2\dd\alpha
 \le CXL^{-A}.
 \tag{1.1}
\]
The order of choice is
\[
 A,\eps\longrightarrow B,D,C,X_0\longrightarrow X,H,J.
\]
The constants may be ineffective.
\end{theorem}

\begin{remark}[Formalization]
The estimate~(1.1) is proved in Lean~4 as
\texttt{AllCenterMAP.map\_two\_fifteenths}.
See \url{https://github.com/jconorgrogan/prime-minor-arcs-2-15};
the formalization has been submitted to Palomar.
\end{remark}

For $h\ne0$ define
\[
 R_X(h)=\sum_{X<n\le2X}\Lambda(n)\Lambda(n+h)
\]
and
\[
 \Sg(h)=
 \begin{cases}
 2C_2\displaystyle\prod_{\substack{p\mid h\\p>2}}\frac{p-1}{p-2},&2\mid h,\\[5pt]
 0,&2\nmid h,
 \end{cases}
 \qquad C_2=\prod_{p>2}\frac{p(p-2)}{(p-1)^2}.
\]

\begin{theorem}[Translated prime-pair consequences]\label{thm:pairs}
Fix $A,\eps>0$.  Uniformly for
\[
 X^{2/15+\eps}\le H\le X^{1-\eps},\qquad
 0\le h_0\le X^{1-\eps},
\]
and $I=\{h\in\Z:|h-h_0|\le H,\ h\ne0\}$,
\[
 \sum_{h\in I}\abs{R_X(h)-X\Sg(h)}^2
 \ll_{A,\eps}HX^2L^{-A},                                      \tag{1.2}
\]
\[
 \sum_{h\in I}R_X(h)^2
 =X^2\sum_{h\in I}\Sg(h)^2+O_{A,\eps}(HX^2L^{-A}),           \tag{1.3}
\]
and
\[
 \#\{h\in I:\abs{R_X(h)-X\Sg(h)}>XL^{-A}\}
 \ll_{A,\eps}HL^{-A}.                                        \tag{1.4}
\]
\end{theorem}

For an integer $N\ge2$, let
\[
 G(N)=\sum_{n+m=N}\Lambda(n)\Lambda(m),
\]
where the pairs are ordered, and let $\Sg(N)$ denote the binary Goldbach
singular series, equal to zero for odd $N$ and to the expression above for even
$N$.

\begin{theorem}[Averaged Goldbach at $2/15$]\label{thm:goldbach}
For every $A,\eps>0$, uniformly for
\[
 X^{2/15+\eps}\le H\le X^{1-\eps},
\]
one has
\[
 G(N)=\Sg(N)N+O_{A,\eps}(XL^{-A})                              \tag{1.5}
\]
for all but $O_{A,\eps}(HL^{-A})$ integers $N\in[X,X+H]$.
The same formula holds for the log-weighted prime-only sum
\[
 W(N)=\sum_{\substack{p+q=N\\p,q\ \mathrm{prime}}}\log p\log q.
\]
Every sufficiently large even target outside the exceptional set
has a Goldbach representation.
\end{theorem}

\section{Two estimates}
The first estimate integrates the two Dirichlet polynomials in \(t\) and in
their relative shift \(u\) at the same time.  The second gives the prime number
theorem simultaneously for every short interval and residue class, for all
moduli \(q\le(\log X)^K\).

\subsection{A two-variable mean for Dirichlet polynomials}

\begin{lemma}[Mixed mean]\label{lem:mixed}
Let $c>0$, let $a\ge0$ be an integer, and let $k\ge1$ be an integer.
Let $M,N\ge2$, $N\ge cM^2$, and $T,U\ge1$.  Suppose
$|\beta(m)|\le\log^a(2m)$ on $(M,2M]$ and
$|g(n)|\le\tau_k(n)\log^a(2n)$ on $(N,2N]$.  Put
\[
 B(t)=\sum_m\beta(m)m^{-1/2-it},\qquad
 G(t)=\sum_ng(n)n^{-1/2-it}.
\]
If $I$ is any interval of length $T$, then, with
$R=\max(2,k^2)$,
\[
 \int_I|G(t)|^2\int_{|u|\le2U}|B(t+u)|^2\dd u\dd t
 \ll_{c,a,k}\log^{4a+2R+2}(2MN)(UT+UN+MN+T).          \tag{2.1}
\]
Fixed character phases may be inserted in either polynomial.
\end{lemma}

\begin{proof}
Use a fixed multiple of
\[
 \eta((t-t_0)/T)\eta(u/U),\qquad
 \eta(v)=\left(\frac{\sin(\pi v/4)}{\pi v/4}\right)^2,
\]
to majorize the rectangle.  The Fourier transform of $\eta$ is supported in
$[-1/4,1/4]$.  Expanding the two squares therefore leaves only quadruples with
\[
 |m_1-m_2|\le K:=\min(M,\pi M/U),\qquad
 |m_1n_1-m_2n_2|\le Y:=2\pi MN/T.
\]
The four weights \(m_i^{-1/2}n_i^{-1/2}\) contribute \(O((MN)^{-1})\), and
the two Fourier transforms contribute \(O(TU)\).  The left side of
(2.1) is at most a fixed power of \(\log(2MN)\) times
\[
 \frac{TU}{MN}
 \sum_{\substack{m_i\asymp M,\ n_i\asymp N\\
 |m_1-m_2|\le K,\ |m_1n_1-m_2n_2|\le Y}}
 \tau_k(n_1)\tau_k(n_2).                               \tag{2.2}
\]
We split this sum into three cases.

If $m_1=m_2$, then $n_1-n_2=O(N/T)$.  Cauchy's inequality and
$\sum_{n\le2N}\tau_k(n)^2\ll N\log^{R-1}(2N)$ give the terms $UT+UN$.

Now let $m_1\ne m_2$.  Write $m_1=da$, $m_2=db$, $(a,b)=1$.
For a nonzero determinant,
\[
 an_1-bn_2=\ell,\qquad0<|\ell|\le Y/d.
\]
Set \(d_1=(\ell,b)\) and \(d_2=(\ell,a)\).  Dividing the first congruence by
\(d_1\) and the second by \(d_2\) gives reduced residue classes, with moduli
\(O_c((N/d_i)^{1/2})\).  Apply Shiu's theorem
\cite[Theorem~1]{Shiu1980} on the full intervals \(N<n_i\le2N\), followed by
Cauchy's inequality.  For each fixed \(\ell\), the result is
\[
 \frac{Nd}{M}\log^{R}(2MN)
 \sqrt{\tau_R(d_1)\tau_R(d_2)}.
\]
Since $(a,b)=1$, $d_1d_2=(\ell,ab)$, while
\[
 \sum_{0<|\ell|\le Y/d}\tau_R((\ell,ab))
 =2\sum_{v\mid ab}\tau_{R-1}(v)\left\lfloor\frac{Y}{dv}\right\rfloor
 \ll\frac Yd\log^{R-1}(2MN).                           \tag{2.3}
\]
If $Y/d<1$, there is no nonzero value of \(\ell\).
Summing the possible pairs \((m_1,m_2)\) gives the term $MN$ in (2.1).

For determinant zero, $n_1=br$ and $n_2=ar$.  Divisor
submultiplicativity and the divisor second moment give
\[
 \sum_r\tau_k(br)\tau_k(ar)
 \ll\frac{Nd}{M}\log^{R-1}(2MN)\tau_k(a)\tau_k(b).
\]
Here $d\le K$.  For fixed $d$, Cauchy's inequality bounds the neighboring
$a,b$ pairs by
\[
 (M/d)(K/d)\log^{R-1}(2MN).
\]
Summing $d\le K$ gives the term $T$.  Combining the three cases proves
(2.1).  Character twists have modulus at most one throughout.
\end{proof}

\subsection{A uniform short-interval prime number theorem}

\begin{proposition}[Simultaneous short-interval progressions]\label{prop:ap}
Fix $D,K>0$ and $\eps>0$.  For all sufficiently large $X$,
\[
 \int_{X/2}^{4X}\max_{q\le L^K}\max_{(a,q)=1}
 \sup_{X^{2/15+\eps}\le Y\le X}
 \left|\frac{\psi(x+Y;q,a)-\psi(x;q,a)-Y/\varphi(q)}Y\right|^2\dd x
 \ll_{D,K,\eps}XL^{-D}.                                \tag{2.4}
\]
The threshold and the implied constant may be ineffective.
\end{proposition}

\begin{proof}
We give the exponent calculation, since it is what fixes $2/15$. It is enough
to treat $0<\eps\le1/10$, since a larger value follows by decreasing $\eps$.
Put
\[
 \theta=\frac2{15}+\eps,\qquad
 \tau=\frac{13}{15}-\frac\eps2,\qquad T=X^\tau,\qquad Q=L^K.
\]
For each character \(\chi\), use the primitive character that induces it.
If \(\rho=\beta+i\gamma\) runs over the nontrivial zeros of the corresponding
\(L\)-function, counted with multiplicity, set
\[
 Z=\sum_{q\le Q}\sum_{\chi\bmod q}\sum_{|\gamma|\le T}X^{2(\beta-1)}.
\]
Set $\delta_0=\min(\eps/16,1/120)$. The density theorem for
polylogarithmic conductors proved in
Appendix~\ref{app:pzd}, used with exponent $K+1$ since
$Q\le(\log T)^{K+1}$ for large $X$, gives, for every fixed
$\eta_{\rm zd}>0$, uniformly on
\[
 \frac12+\delta_0\le\sigma\le\frac45,
\]
\[
 \sum_{q\le Q}\sum_{\chi\bmod q}N(\sigma,T,\chi)
 \ll_{K,\delta_0,\eta_{\rm zd}}
 T^{\frac{30}{13}(1-\sigma)+\eta_{\rm zd}}.          \tag{2.5}
\]
Its inputs are Montgomery's classical Dirichlet-$L$ zero detector and moment
estimates together with Guth--Maynard's published arbitrary-coefficient
large-values theorem, applied separately to each of the only polylogarithmically
many characters.  Thus no hybrid character-family theorem is used. The zeros with
$\beta<1/2+\delta_0$ contribute
\[
 \ll Q^2T L X^{-1+2\delta_0}\ll_{\eps,K}X^{-c_\eps}
\]
by the elementary zero count.  Partition
$[1/2+\delta_0,4/5]$ into $O_\eps(1)$ intervals
$[\sigma,\sigma+\Delta]$ with
\[
 \Delta\le\min\!\left(\frac1{100},\frac{3\eps}{52}\right),
\]
and apply (2.5) at each left endpoint. Choose
\[
 \eta_{\rm zd}=\frac{3\eps}{52\tau}.
\]
Since
\[
 2-\frac{30}{13}\tau=\frac{15}{13}\eps>0,
\]
the $X$-exponent of a mesh cell is at most
\[
 -\frac{15\eps}{13}(1-\sigma)+2\Delta+\tau\eta_{\rm zd}
 \le-\frac{3\eps}{52}.
\]
Here we used $\sigma\le4/5$ and
$\Delta\le3\eps/52$. Thus every mesh interval contributes a fixed negative
power of $X$, uniformly over all $q\le Q$.

Near one, Jutila \cite[(1.7)]{Jutila1977} proves, for every fixed $\eta>0$,
\[
 \sum_{\chi\bmod q}N(\sigma,T,\chi)
 \ll_\eta(qT)^{(2+\eta)(1-\sigma)},\qquad4/5\le\sigma<1. \tag{2.6}
\]
With $\eta=1/10$,
\[
 2-\frac{21}{10}\frac{\log(qT)}{\log X}\ge\frac16
\]
for large $X$.  Khale's Vinogradov--Korobov region
\cite[Theorem~1.1]{Khale2024} gives, for every zero with $|\gamma|\ge10$,
\[
 1-\beta\gg_K L^{-2/3}(\log L)^{-1/3}.
\]
For the conductor-one zeta factor, apply Khale's theorem to the principal
character modulo~$3$; the additional factor $1-3^{-s}$ has zeros only on
$\Re s=0$.
The classical low-height region and Siegel's theorem handle the
possible real exceptional zero; for $q\le L^K$ its weighted contribution is
$O(L^{2K}\exp(-c_KL^{3/4}))$.  For every fixed $D$,
\[
 Z\ll_{D,K,\eps}L^{-D}.                                 \tag{2.7}
\]
The explicit formula and local zero count used here are
\cite[Theorem~11.3, Lemma~11.4(a), Theorems~12.3 and 12.10]{Koukoulopoulos}.
They apply to arbitrary characters; changing to the inducing character costs
$O(\log q\log X)$.

Let
\[
 F_\chi(t)=\sum_{|\gamma|\le T}t^{\rho-1}
 \qquad (X/4\le t\le6X),
\]
extended by zero, where the zeros are those of the primitive character
inducing~$\chi$.  If $\rho=\beta+i\gamma$ and
$\rho'=\beta'+i\gamma'$, direct integration gives
\[
 \left|\int_{X/4}^{6X}t^{\rho+\overline{\rho'}-2}\,\dd t\right|
 \ll \frac{X^{\beta+\beta'-1}}{1+|\gamma-\gamma'|}.
\]
Writing $w_\rho=X^{\beta-1}$, using
$2w_\rho w_{\rho'}\le w_\rho^2+w_{\rho'}^2$, and grouping the ordinates
of~$\rho'$ into unit intervals, Lemma~11.4(a) and the harmonic sum give
\[
 \norm{F_\chi}_2^2
 \ll XL^2\sum_{|\gamma|\le T}X^{2(\beta-1)}.
\]
Consequently
\[
 \sum_{q\le Q}\sum_{\chi\bmod q}\norm{F_\chi}_2^2
 \ll XL^2Z.                                             \tag{2.8}
\]

Let
\[
 E_\chi(x,Y)=\psi(x+Y,\chi)-\psi(x,\chi)
              -\mathbf 1_{\chi=\chi_0}Y.
\]
Subtracting the explicit formula at $x+Y$ and $x$, and using the
primitive--imprimitive correction, gives uniformly in the asserted ranges
\[
 \frac{|E_\chi(x,Y)|}{Y}
 \le M^+F_\chi(x)+O(X^{-\eps/2}L^2),                    \tag{2.9}
\]
where $M^+$ is the one-sided Hardy--Littlewood maximal operator.  Character
orthogonality is used only before the supremum:
\[
 E_{q,a}(x,Y)=\frac1{\varphi(q)}
   \sum_{\chi\bmod q}\overline\chi(a)E_\chi(x,Y).
\]
After (2.9), Cauchy--Schwarz, rather than orthogonality, yields
\[
 \sup_Y\left|\frac{E_{q,a}(x,Y)}Y\right|^2
 \le \frac1{\varphi(q)}\sum_{\chi\bmod q}
   \left(M^+F_\chi(x)+O(X^{-\eps/2}L^2)\right)^2.       \tag{2.10}
\]
The right side is independent of~$a$.  Since the summands are nonnegative,
$\max_{q\le Q}$ is at most their sum over~$q$.  The $L^2$ boundedness of
$M^+$, followed by (2.7)--(2.8), therefore gives
\[
 \int_{X/2}^{4X}\max_{q\le Q}\max_{(a,q)=1}\sup_Y
 \left|\frac{E_{q,a}(x,Y)}Y\right|^2\,\dd x
 \ll XL^2Z+QX^{1-\eps}L^4.
\]
Taking the logarithmic saving in (2.7) larger by two, and absorbing the
second, power-saving term, proves (2.4).
\end{proof}

\begin{remark}[Comparison with the hybrid theorem]
Chen--Gupta--Li version~2 \cite[Theorem~1.2]{CGL2026} states the stronger
hybrid estimate
\[
 \sum_{\chi\bmod q}N(\sigma,T,\chi)
 \le(qT)^{o(1)}\left(q^{\frac73(1-\sigma)}T^{2(1-\sigma)}
 +(qT)^{\frac{30}{13}(1-\sigma)}\right).
\]
That theorem also implies Proposition~\ref{prop:ap}, but it is not used in
this paper. Appendix~\ref{app:pzd} proves the weaker polylogarithmic-conductor
slice needed here from published inputs.
\end{remark}

\section{Near and far from rational points}
Write the center of a short Fourier interval as
\(\gamma=a/q+\lambda\).  If \(\lambda\) is small,
Proposition~\ref{prop:ap} controls a neighborhood of \(a/q\).  If
\(\lambda\) is larger, Matom\"aki--Radziwi\l\l--Tao convert the Fourier
integral into the two-variable mean from Lemma~\ref{lem:mixed}.  We split at
\(|\lambda|H_*\asymp R\); both estimates use the same values of \(a\) and
\(q\).

\begin{proof}[Proof of Theorem~\ref{thm:map}]
Put
\[
 \rho=\min(\eps/4,1/1200),\quad
 H_*=\tfrac12X^{2/15+\rho},\quad
 \delta=\rho/16.
\]
Every arc of length $1/H$, $H\ge X^{2/15+\eps}$, lies in the concentric arc
of length $1/H_*$.  It is therefore enough to work at $H_*$.  Set
\[
 K_0=\lceil1/\delta\rceil,\quad \kappa=2K_0,
\]
choose \(C_s=C_s(\eps)\) large enough to cover the fixed powers of \(\log X\)
in the Heath--Brown decomposition and its mean-value estimates, and then take
\[
 B=\left\lceil4(A+C_s+20)\right\rceil,\quad Q=L^B,
\]
\[
 C_c=A+C_s+2B+30,\quad R=L^{C_c},\quad
 y=\frac{H_*}{32R},
\]
\[
 D_{\rm AP}=A+4B+10,\qquad
 D=\left\lceil A+2B+10\right\rceil,\qquad P=L^D.     \tag{3.1}
\]
These numbers are fixed before $X$, the arc length, or the center is chosen.
Apply Proposition~\ref{prop:ap} with
$\eps_{\rm AP}=\rho/2$, $K=B$, and $D=D_{\rm AP}$.  After enlarging the
lower threshold for $X$, we have
$64L^{C_c}\le X^{\rho/2}$, and hence
$y\ge X^{2/15+\rho/2}$.  Thus the proposition applies at length $y$ with
exactly the denominator range and logarithmic saving used below.

\paragraph{Near a rational point.}
For each reduced $a/q$, introduce the signed measure
\[
 \nu_{q,a}=\sum_{X<n\le2X}\Lambda(n)e(an/q)\delta_n
 -\frac{\mu(q)}{\varphi(q)}\mathbf1_{(X,2X]}(t)\,dt.
\]
Since $4R/H_*=1/(8y)$, the signed-measure Gallagher inequality gives
\[
 \int_{|\beta|\le1/(8y)}|\widehat\nu_{q,a}(\beta)|^2\,d\beta
 \ll y^{-2}\int_{\mathbb R}|\nu_{q,a}((x,x+y])|^2\,dx.
\]
Decompose the interior sum into reduced residue classes modulo $q$ and apply
Cauchy's inequality before Proposition~\ref{prop:ap}.  The interior,
endpoint, and prime-power terms contribute respectively
$q^2XL^{-D_{\rm AP}}$, $yL^2$, and $L^4/y$ for one rational arc after
Gallagher normalization.  Summing the reduced numerators over $q\le Q$ and
using $\sum_{q\le Q}\varphi(q)q^2\ll Q^4$ and
$\sum_{q\le Q}\varphi(q)\ll Q^2$, while estimating the continuous main term
outside $|\beta|\le P/X$ by $|\beta|^{-1}$, gives on the union $\mathcal W$ of intervals
$\|\alpha-a/q\|\le4R/H_*$, $q\le Q$,
\[
 \int_{\mathcal W\setminus\M(P,Q)}|S_X(\alpha)|^2\dd\alpha
 \ll XQ^4L^{-D_{\rm AP}}+Q^2yL^2+\frac{Q^2L^4}{y}+\frac{XQ^2}{P}
 \ll XL^{-A}.                                          \tag{3.2}
\]
The last term comes from the continuous main term
$\int_X^{2X}e(\beta t)\dd t$ outside $|\beta|\le P/X$; prime powers sharing
a factor with $q$ and the two endpoints of $(X,2X]$ give the middle terms.
Every term uses the same
denominator range \(q\le Q\).

\paragraph{Farther from a rational point.}
Let $\gamma$ be the midpoint of a target arc.  Dirichlet approximation gives
\[
 \gamma=a/q+\lambda,\qquad q\le Q,\qquad |\lambda|\le(qQ)^{-1}.
\]
If $|\lambda|\le2R/H_*$, the target interval lies in $\mathcal W$, and (3.2)
applies.  Otherwise put $U=|\lambda|H_*>2R$.  Corollary~5.3 of
Matom\"aki--Radziwi\l\l--Tao \cite{MRT2019}, with $\eta=Q^{-1/2}$, gives
\[
 \int_{\lambda-1/H_*}^{\lambda+1/H_*}|S_X(a/q+\theta)|^2\dd\theta
 \ll\frac{d_2(q)^4}{qU^2}\sup_{q=q_0q_1}\mathcal I(q_0,q_1)+\mathcal E_0,
 \tag{3.3}
\]
where
\[
 \mathcal I(q_0,q_1)=\int_I\left(
 \sum_{\chi\bmod q_1}\int_{t-U}^{t+U}
 |\mathcal D[\Lambda\1_{(X,2X]}](1/2+it',\chi,q_0)|\dd t'
 \right)^2\dd t
\]
and \(I\) is the range
$Q^{-1/2}|\lambda|X\le|t|\le Q^{1/2}|\lambda|X$.
The denominator in (3.3) is $qU^2$.  The remaining term satisfies
$\mathcal E_0\ll XL^{O(1)}(Q^{-1/2}+U^{-1})^2$, which is
$O(XL^{-A})$ by (3.1).

Use the $m=8$ Heath--Brown decomposition in
\cite[Lemmas~2.15--2.16]{MRT2019}.  Its hypotheses hold because, with
$H_0=H_*X^{-\delta}$,
\[
 (2/15+\rho-\delta)-(1/8+\delta)=1/120+7\rho/8>0.       \tag{3.4}
\]
The small terms save a fixed power of $X$.  The terms with $q_0>1$ are
supported on powers of the prime dividing $q_0$ and have squared norm
$O(L^3)$; hence \(q_0=1\) in the remaining convolution terms.

Remove the sharp product cutoff with
\cite[Corollary~2.5]{MRT2019}, applied to the \emph{untruncated}
convolution $f'$.  Factor the resulting polynomial and then restore the
cutoff.  The Perron error saves a fixed power of \(X\), and its translations
stay inside a fixed enlargement of the same \(t\)-range.

\begin{remark}[Notation in the truncation-removal step]
In the corresponding display of \cite[Section~6]{MRT2019}, the occurrences
of $\widetilde f$ in the application of Corollary~2.5 and in the definition
of $F$ appear to be typographical slips for the untruncated convolution
$f'$.  This interpretation is forced by the preceding identity,
Proposition~6.1, and the parallel Appendix~A argument.  We write the
untruncated function explicitly and pay the truncation error above.
\end{remark}

The decomposition labels its remaining pieces Type II and Type \(d_j\).
Type II and the first two Type-\(d\) pieces follow from the character second-
and fourth-moment estimates in \cite[Lemma~2.10 and Corollary~2.12]{MRT2019},
as in Section~7 there.  After
division by \(qU^2X\), the resulting terms are
\[
 qQ^{1/2}|\lambda|,\quad Q^{1/2}X^{-\delta},\quad
 X^{-\delta},\quad U^{-1},                              \tag{3.5}
\]
all small in the present range.

For a Type-$d_j$ cell, $3\le j\le7$, select its shortest displayed
$\beta$-factor, of length $M$.  At least two other displayed $\beta$-factors
are no shorter, so the
complementary length $N'$ satisfies $N'\gg M^2$.  Dyadically decompose the
complementary convolution into $O_\eps(1)$ blocks of length $N'$.  On every
block, $N'\gg NM_2\cdots M_j\gg M^2$, and its coefficient remains bounded by
a fixed divisor function, at the cost of a fixed logarithmic factor.  Cauchy
and Fubini bound the overlap of the two inner $t'$-windows by $O(U)$;
expanding the squared character sum leaves at most $q^2$ mixed means.
Lemma~\ref{lem:mixed}
then gives, after division by $qU^2X$,
\[
 L^{O(1)}\left(qQ^{1/2}|\lambda|+\frac qM+\frac qU
 +\frac{qQ^{1/2}}{H_*}\right)                           \tag{3.6}
\]
which is at most
\[
 L^{O(1)}\left(Q^{-1/2}+QX^{-2/15-\rho+\delta}
 +Q/R+Q^{3/2}/H_*\right).                               \tag{3.7}
\]
The first and third terms save powers of \(L\); the other terms save powers of
\(X\).  With the values in (3.1), multiplying each term by \(L^{C_s}\)
still gives \(O(L^{-A})\).  The definition of \(C_s\) includes the number of
pieces, both character sums, the Perron shifts, and the divisor factor in
(3.3).

The two cases cover every center.  In the second case, (3.3) bounds the entire
target interval and therefore its minor-arc part.  Together with (3.2), this
gives (1.1).
\end{proof}

\begin{remark}[The major arcs must depend on the saving]
A rational peak at a denominator $r$ just beyond a fixed cutoff $Q$ has local
energy of order $X/r^2$, which exceeds $XL^{-A}$ once $A>2B+O(1)$, and the
tail of the continuous main term forces the width exponent to grow with $A$
as well.  The order $\forall A\,\exists B,D$ in Theorem~\ref{thm:map} is
therefore essential; the reversed order $\exists B,D\,\forall A$ gives a
false statement.
\end{remark}

\section{Prime-pair variance, fourth correlation, and density one}
\begin{proof}[Proof of Theorem~\ref{thm:pairs}]
Let $f(n)=\Lambda(n)\1_{(X,2X]}(n)$ and
$C_f(h)=\sum_nf(n)f(n+h)$.  Use the set
\(\M=\M(P,Q)\) supplied by Theorem~\ref{thm:map}, put
\(\m=\T\setminus\M\), and define the major-arc contribution
\[
 M_f(h)=\int_{\M}|S_X(\alpha)|^2e(-h\alpha)\dd\alpha.
\]
A Fejér-kernel large sieve for the positive measure
$\1_{\m}|S_X(\alpha)|^2\dd\alpha$ gives
\[
 \sum_{|h-h_0|\le H}|C_f(h)-M_f(h)|^2
 \ll H\left(\int_{\m}|S_X|^2\right)
 \sup_{|J|=1/H}\int_{J\cap\m}|S_X|^2.                  \tag{4.1}
\]
Parseval gives $\int_\T|S_X|^2=\sum_{X<n\le2X}\Lambda(n)^2\ll XL$.
Use Theorem~\ref{thm:map} with one extra logarithm to obtain
\[
 \sum_{|h-h_0|\le H}|C_f(h)-M_f(h)|^2\ll HX^2L^{-A}.  \tag{4.2}
\]

The standard major-arc calculation, using
\cite[Proposition~4.1]{MRT2019}, yields, for $0<|h|<X$,
\[
 M_f(h)=(X-|h|)\Sg(h)+O_C(XL^{-C}\tau(|h|)^2).          \tag{4.3}
\]
Standard mean-value bounds for fixed powers of \(\tau\) control the sum of
these errors after $C$ is enlarged.  The difference between $C_f(h)$ and
$R_X(h)$, together with
the overlap correction $|h|\Sg(h)$, is
$O((|h|+1)L^{O(1)})$ per shift.  Since the stated range gives
$|h|=O(X^{1-\eps})$, its total square is smaller by a fixed power of \(X\)
than the right side of (1.2).  Adding these contributions to (4.2) gives
(1.2).

Write $E_h=R_X(h)-X\Sg(h)$.  Then
\[
 \sum_{h\in I}R_X(h)^2
 =X^2\sum_{h\in I}\Sg(h)^2
 +2X\sum_{h\in I}\Sg(h)E_h+\sum_{h\in I}E_h^2.        \tag{4.4}
\]
The standard bound $\sum_{h\in I}\Sg(h)^2\ll HL^{O(1)}$, Cauchy's inequality,
and (1.2) with a larger logarithmic exponent make the last two terms
$O(HX^2L^{-A})$, giving (1.3).

Apply (1.2) with exponent $3A$.  Every exception counted in (1.4)
contributes more than $X^2L^{-2A}$ to the variance, which gives (1.4).
\end{proof}

\section{Averaged Goldbach}
\begin{proof}[Proof of Theorem~\ref{thm:goldbach}]
Choose every cutoff independently of the target.  For integers $r\ge1$ define
\[
 L_X(r)=\1_{X/2<r\le X},\qquad Q_X(r)=\1_{1\le r<X},
\]
\[
 A_X(N)=\sum_{n+m=N}\Lambda(n)\Lambda(m)L_X(n)Q_X(m),
\quad
 B_X(N)=\sum_{n+m=N}\Lambda(n)\Lambda(m)L_X(n)L_X(m),
\]
and
\[
 E_X(N)=\sum_{n+m=N}\Lambda(n)\Lambda(m)
 \{1-L_X(n)Q_X(m)-Q_X(n)L_X(m)+L_X(n)L_X(m)\}.
\]
The pointwise indicator identity gives
\[
 G(N)=2A_X(N)-B_X(N)+E_X(N).                            \tag{5.1}
\]
Because \(Q_X\) excludes \(X\), the overlap in \(A_X\) excludes \(X\) as
well.  The endpoint term is
\[
 2\,\1_{X/2<N-X<X}\Lambda(X)\Lambda(N-X)
\]
when $X$ is an integer.  For sufficiently large $X$, however,
$H\le X^{1-\eps}\le X/2$, so this term vanishes.  Writing $N=X+h$, the sum
\(B_X(N)\) and the terms with a summand greater than \(X\) each contain
$O(h+1)$ pairs.  Therefore
\[
 |B_X(N)|+|E_X(N)|\ll(H+1)L^2\ll_{A,\eps}XL^{-A}.       \tag{5.2}
\]

Put
\[
 S(\alpha)=\sum_{X/2<n\le X}\Lambda(n)e(n\alpha),\qquad
 F(\alpha)=\sum_{1\le m<X}\Lambda(m)e(m\alpha).
\]
Orthogonality gives
\[
 A_X(N)=\int_\T S(\alpha)F(\alpha)e(-N\alpha)\dd\alpha.\tag{5.3}
\]
Let \(\M\) be the union of the rational intervals used in the major-arc
calculation below, and put \(\m=\T\setminus\M\).  On \(\m\), define
$d\nu=\1_\m SF\dd\alpha$.  For any finite complex measure $\nu$ and any
block $I$ of at most $H+1$ consecutive integers, the nonnegative Fejér kernel
gives
\[
 \sum_{k\in I}|\widehat\nu(k)|^2
 \ll H|\nu|(\T)\sup_{|J|=1/H}|\nu|(J).                \tag{5.4}
\]
To prove (5.4), shift the frequency block to the origin and use the Fej\'er
kernel $K_{2\lceil H\rceil+2}$, whose triangular Fourier weights are bounded
below on that block, including its endpoints.  After taking total variation,
partition the tail using
$K_{2\lceil H\rceil+2}(t)\ll H(1+H\|t\|_\T)^{-2}$.

Parseval and Cauchy give $|\nu|(\T)\ll XL$.  Apply
Theorem~\ref{thm:map} to $S$ at scale $X/2$ with saving $K$.  Replacing
$\eps$ by $\eps/2$ makes both bounds
\[
 (X/2)^{2/15+\eps/2}\le H\le(X/2)^{1-\eps/2}
\]
hold for large $X$.  Another Cauchy inequality gives
\[
 \sup_{|J|=1/H}|\nu|(J)\ll XL^{(1-K)/2}.
\]
\[
 \sum_{X\le N\le X+H}|\widehat\nu(N)|^2
 \ll HX^2L^{(3-K)/2}.                                  \tag{5.5}
\]
At the pointwise threshold $XL^{-A}$, Chebyshev's inequality gives
$O(HL^{-A})$ exceptions when
\[
 K\ge6A+3.                                             \tag{5.6}
\]

For the major arcs, let $Q=L^B$ and use width
$L^{B_0}/X$.  Choose $B,B_0$ after the exponents defining the MAP arcs at
scale $X/2$, and large enough that the present major-arc set contains those
arcs.  For any fixed \(R>0\), Proposition~4.1 of \cite{MRT2019}
gives, when $\alpha=a/q+\beta$,
\[
 S(\alpha)=\frac{\mu(q)}{\varphi(q)}\int_{X/2}^{X}e(\beta x)\dd x+O(XL^{-R}),
\]
\[
 F(\alpha)=\frac{\mu(q)}{\varphi(q)}\int_0^{X}e(\beta y)\dd y+O(XL^{-R}).
\]
For $X\le N\le X+H\le3X/2$, the continuous convolution has the exact length
\[
 \int_{\mathbb R}\1_{(X/2,X]}(x)\1_{[0,X)}(N-x)\dd x=X/2. \tag{5.7}
\]
After the $\beta$ integral and the reduced-residue sum,
\[
 \int_\M SF e(-N\alpha)\dd\alpha
 =\frac X2\Sg(N)+O\!\left(XL^{-C_{\rm major}}\tau(N)^{O(1)}\right). \tag{5.8}
\]
This derivation uses Proposition~4.1 of \cite{MRT2019} directly; the
displayed Goldbach major-arc proposition there has supports whose continuous
overlap vanishes at the displayed coefficient, so it is bypassed.
The condition $B\ge2C_{\rm major}$ bounds the omitted part of the Ramanujan
series, and $B_0\ge2B+C_{\rm major}$ bounds the extension of the \(\beta\)
integral.  A short-interval mean-value bound for a fixed power of \(\tau\),
followed by Markov's inequality, gives at most
$O(HL^{-A})$ targets on which the error in (5.8) exceeds $XL^{-A}$.

Choose the parameters in the order
\[
 A,\eps\to C_{\rm major}\to K\to B\to B_0,R\to X_0.
\]
Outside the two exceptional sets, (5.3), (5.5), and (5.8) give
\[
 A_X(N)=\frac X2\Sg(N)+O_{A,\eps}(XL^{-A}).
\]
Equations (5.1)--(5.2) yield $G(N)=X\Sg(N)+O(XL^{-A})$.
Also,
$|(N-X)\Sg(N)|\ll H L^{O(1)}\ll XL^{-A}$, after enlarging the input saving,
so (1.5) follows.

Every term of $G(N)-W(N)$ contains a proper prime power.  Chebyshev's bound
gives uniformly
\[
 |G(N)-W(N)|\ll X^{1/2}L^2,
\]
smaller than $XL^{-A}$ for every fixed $A$.  The weighted prime-only formula
follows.  For even \(N\) outside the exceptional set, \(\Sg(N)\) is
bounded below by a positive constant, so $W(N)>0$ when $X$ is large.
\end{proof}

There is a balanced variant, with both summands in the same dyadic range.

\begin{theorem}[Balanced averaged Goldbach]\label{thm:balanced}
Fix $A,\eps>0$ and $\delta\in(0,1/2)$.  For an integer $N\in(2X,4X)$ put
\[
 G_X(N)=\sum_{\substack{n+m=N\\X<n,m\le2X}}\Lambda(n)\Lambda(m),
 \qquad
 \ell_X(N)=\min(N-2X,\,4X-N).
\]
Uniformly for $X^{2/15+\eps}\le H\le X^{1-\eps}$ and every interval
$J\subset[(2+\delta)X,(4-\delta)X]$ of length $H$,
\[
 G_X(N)=\Sg(N)\,\ell_X(N)+O_{A,\eps,\delta}(XL^{-A})
 \tag{5.9}
\]
for all but $O_{A,\eps,\delta}(HL^{-A})$ integers $N\in J$.  Outside the
same exceptional set, every even $N\in J$ satisfies
\[
 \#\{p+q=N:X<p,q\le2X\}
 =\frac{\Sg(N)\,\ell_X(N)}{L^2}\left(1+O\!\left(\frac1L\right)\right).
\]
\end{theorem}

\begin{proof}
Both factors are now $S_X$, so the argument above simplifies.  Orthogonality
gives $G_X(N)=\int_\T S_X(\alpha)^2e(-N\alpha)\dd\alpha$.  Take $\M$ as in
the major-arc calculation above and put $d\nu=\1_\m S_X^2\dd\alpha$.
Parseval gives $|\nu|(\T)\le\sum_{X<n\le2X}\Lambda(n)^2\ll XL$, while
Theorem~\ref{thm:map} with saving $K$ gives directly
\[
 \sup_{|J'|=1/H}|\nu|(J')
 \le\sup_{|J'|=1/H}\int_{J'\cap\m}|S_X(\alpha)|^2\dd\alpha
 \ll XL^{-K}.
\]
No rescaling is needed, because $S_X$ is the sum in Theorem~\ref{thm:map}
itself.  Now (5.4) gives
\[
 \sum_{N\in J}|\widehat\nu(N)|^2\ll HX^2L^{1-K}.
\]
Chebyshev's inequality at the threshold $XL^{-A}$ gives $O(HL^{-A})$
exceptions when $K\ge3A+1$.
On the major arcs, the approximation of $S_X$ by
Proposition~4.1 of \cite{MRT2019}, the $\beta$ integral, and the
reduced-residue sum give, exactly as in (5.8),
\[
 \begin{aligned}
 \int_\M S_X(\alpha)^2e(-N\alpha)\dd\alpha
 &=\Sg(N)\int_{\mathbb R}\1_{(X,2X]}(x)\1_{(X,2X]}(N-x)\dd x\\
 &\quad+O\!\left(XL^{-C_{\rm major}}\tau(N)^{O(1)}\right),
 \end{aligned}
\]
and the continuous overlap equals $\ell_X(N)$.  The divisor-moment and
Markov step is unchanged.  Prime powers contribute $O(X^{1/2}L^2)$
uniformly.  On $J$ one has $\ell_X(N)\ge\delta X$, and every surviving
summand satisfies $\log p\log q=L^2(1+O(1/L))$, which gives the prime-only
count.
\end{proof}

\section*{Scope and machine-checked components}
The results concern almost every shift or target in every sufficiently long
translated interval.  A prescribed shift such as $h=2$, a prescribed
Goldbach target, and a pair of prescribed short summand intervals are
pointwise problems; the present methods leave them open.

Siegel's theorem enters Proposition~\ref{prop:ap} through the possible
exceptional zero, so the constants throughout may be ineffective.

Theorem~\ref{thm:map} is formalized in Lean~4 as
\texttt{AllCenterMAP.map\_two\_fifteenths}.  The proof is
\url{https://github.com/jconorgrogan/prime-minor-arcs-2-15}
and has been submitted to Palomar, with
\texttt{comparator.json} targeting that declaration.  The release passed a
Linux source build of $10{,}255$ jobs, Comparator statement identity against
the independently compiled Challenge, independent NanoDa replay, and Lean
kernel acceptance.  The MAP proof closure contains $100{,}596$ declarations
across $1{,}861$ modules; the internal Guth--Maynard large-value closure
contains $75{,}874$ declarations.  Both closures replayed into an empty
kernel.  The axiom audit reports only \texttt{propext},
\texttt{Classical.choice}, and \texttt{Quot.sound}.

\appendix
\section{A published-input density lemma for polylogarithmic conductors}
\label{app:pzd}

The all-center argument only uses conductors bounded by a fixed power of the
logarithm. The following narrower lemma avoids any appeal to an unrefereed
hybrid character-aspect theorem while retaining the $30/13$ coefficient.

For a Dirichlet character $\chi$, let $N(\sigma,T,\chi)$ count, with
multiplicity, the zeros $\rho=\beta+i\gamma$ of $L(s,\chi)$ satisfying
$\beta\ge\sigma$ and $|\gamma|\le T$.

\begin{theorem}[Polylogarithmic-conductor density]\label{thm:pzd}
Fix $K>0$, $\delta>0$, and $\eta>0$.  Uniformly for
\[
 T\ge T_0(K,\delta,\eta),\qquad
 \frac12+\delta\le\sigma\le\frac45,
 \qquad Q\le(\log T)^K,
\]
one has
\[
 \sum_{q\le Q}\sum_{\chi\bmod q}N(\sigma,T,\chi)
 \ll_{K,\delta,\eta}
 T^{\frac{30}{13}(1-\sigma)+\eta}.
 \tag{A.1}
\]
\end{theorem}

\begin{proof}
We first reduce to primitive characters.  If $\chi\pmod q$ is induced by the
primitive character $\chi^*\pmod r$, then
\[
 L(s,\chi)=L(s,\chi^*)
 \prod_{\substack{p\mid q\\p\nmid r}}
       (1-\chi^*(p)p^{-s}).
\]
Every zero of a factor in the finite product has real part zero.  Consequently
the zeros counted in (A.1) are zeros of $L(s,\chi^*)$.  The number of
character--modulus pairs under consideration is $O(Q^2)$, so it is enough to
prove, for every primitive $\chi\pmod r$, $r\le Q$,
\[
 N(\sigma,T,\chi)
 \ll_{K,\delta,\eta}
 T^{\frac{30}{13}(1-\sigma)+\eta/2}.
 \tag{A.2}
\]

We give the reduction of (A.2) to published large-values and moment estimates.
This also records where conductor uniformity enters. Fix small constants
$\kappa>0$ and $0<\vartheta<1$, to be chosen at the end in terms of $\eta$, and put
\[
 \mathcal R=r(T+2),\qquad U=2\mathcal R^\kappa,
 \qquad Y=\mathcal R^{1/2}.
\]
If $r=1$, the character is the zeta character.  The range
$7/10\le\sigma\le4/5$ follows from \cite[Theorem~1.2]{GuthMaynard2026},
while the lower range follows from the Ingham estimate summarized in
\cite[(1.4)]{GuthMaynard2026}.  We may therefore assume $r>1$;
the primitive character $\chi$ is then nonprincipal and no pole residue
occurs. Form the usual mollifier
\[
 M_U(s,\chi)=\sum_{d\le U}\mu(d)\chi(d)d^{-s}.
\]
For $\operatorname{Re} s>1$,
\[
 L(s,\chi)M_U(s,\chi)
 =\sum_{n\ge1}c_n\chi(n)n^{-s},
 \qquad c_n=\sum_{\substack{d\mid n\\d\le U}}\mu(d).
 \tag{A.3}
\]
Thus $c_1=1$ and $c_n=0$ for $1<n\le U$. We now record the detector without
discarding its weights. Mellin inversion for $e^{-n/Y}$ and a contour shift
to $\operatorname{Re} z=1/2-\beta$ give, at a zero
$\rho=\beta+i\gamma$,
\begin{align}
 e^{-1/Y}+\sum_{n>U}c_n\chi(n)n^{-\rho}e^{-n/Y}
  ={}&\frac1{2\pi i}\int_{\operatorname{Re} z=1/2-\beta}
       \Gamma(z)Y^zL(\rho+z,\chi)M_U(\rho+z,\chi)\,dz.         \tag{A.4}
\end{align}
The apparent pole at $z=0$ is removable: since $L(\rho,\chi)=0$,
$L(\rho+z,\chi)=zH_\rho(z)$ with $H_\rho$ holomorphic near zero, while
$z\Gamma(z)=\Gamma(z+1)$.  There is no pole at $z=1-\rho$ because $\chi$
is nonprincipal.  If
$N_\chi(t,t+1)$ denotes the number of nontrivial zeros in the critical strip,
counted with multiplicity, whose ordinates lie in $[t,t+1]$, we use the local
zero count
\[
 N_\chi(t,t+1)\ll\log(r(|t|+2)).                                \tag{A.5}
\]
This estimate controls multiplicity.  Put
$a=1/2-\beta$ and $B=(\log\mathcal R)^2$.  Parametrizing the shifted line by
$z=a+iv$ gives $dz=i\,dv$ and $\rho+z=1/2+i(\gamma+v)$.  Equivalently, with
the absolute critical-line ordinate $u=\gamma+v$, the integrand is
\[
 \Gamma\bigl(a+i(u-\gamma)\bigr)
 Y^{a+i(u-\gamma)}L(1/2+iu,\chi)M_U(1/2+iu,\chi),
\]
and the truncated interval is $|u-\gamma|\le B$.  Uniformly for
$1/2<\beta\le1$ and $|v|\ge1$, the fixed-strip bound for $L$, the elementary
bound $|M_U(1/2+i(\gamma+v),\chi)|\le U+1$, and vertical Gamma decay give
\begin{align*}
 &\left|\Gamma(a+iv)Y^{a+iv}
 L(1/2+i(\gamma+v),\chi)M_U(1/2+i(\gamma+v),\chi)\right| \\
 &\qquad\ll r^2(U+1)(1+|v|)(5+|\gamma|+|v|)^2e^{-|v|} \\
 &\qquad\ll r^2(U+1)(5+|\gamma|)^2(1+|v|)^3e^{-|v|}.
\end{align*}
Hence the two vertical tails are
\[
 \ll r^2(U+1)(5+|\gamma|)^2e^{-B/2}.
\]
If $N=\lceil YB\rceil$, the arithmetic tail is at most
\[
 (U+1)e^{-(N+1)/Y}(1-e^{-1/Y})^{-1}
 \ll (U+1)Y e^{-B}.
\]
Since $\mathcal R=r(T+2)$, $Y=\mathcal R^{1/2}$,
$U=2\mathcal R^\kappa$, $|\gamma|\le T$, and
$r\le(\log T)^K$, both tails are $O_A(\mathcal R^{-A})$ for every fixed
$A$. Consequently every remaining zero satisfies either
\[
 \left|\sum_{U<n\le Y(\log\mathcal R)^2}
 c_n\chi(n)n^{-\rho}e^{-n/Y}\right|\gg1,                        \tag{I}
\]
or the shifted critical-line integral containing
$L(1/2+i(\gamma+u),\chi)M_U(1/2+i(\gamma+u),\chi)$ is large for
$|u|\le(\log\mathcal R)^2$. This is the classical Montgomery
Dirichlet-$L$ detector \cite[Chapters~10 and 12]{Montgomery1971}; (A.3) displays
the character-twisted coefficient identity. Equation (A.5) controls
multiplicity and permits thinning at any prescribed separation, with the
separation length times a logarithm as the loss.

For the critical-line alternative, first thin the original ordinates at
separation $3(\log\mathcal R)^2$. The shifted ordinates are then
one-separated, and the thinning loss is polylogarithmic by (A.5). The detector
inequality and the integrability of the gamma factor imply, for one shifted
ordinate $v$, that
\[
 |L(1/2+iv,\chi)M_U(1/2+iv,\chi)|
 \gg Y^{\sigma-1/2}\mathcal R^{-o(1)}.
\]
Since
\[
 |M_U(1/2+iv,\chi)|\le\sum_{d\le U}d^{-1/2}\ll U^{1/2},
\]
Montgomery's discrete family fourth moment
\cite[Theorem~10.3]{Montgomery1971}, restricted to the fixed primitive character
by nonnegativity,
\[
 \sum_v|L(1/2+iv,\chi)|^4\ll(rT)^{1+o(1)}
\]
at one-separated ordinates gives
\[
 \#\mathcal Z_{II}\ll\mathcal R^{2(1-\sigma)+2\kappa+o(1)}.           \tag{A.6}
\]
All conductor factors are harmless because
$r\le(\log T)^K$. Notice that we used the pointwise bound for $M_U$; a
fourth moment of $L$ and only a second moment of $M_U$ would not justify (A.6).

It remains to count the zeros satisfying (I). Dyadic subdivision supplies
one $N$ with
\[
 \mathcal R^\kappa\ll N\ll\mathcal R^{1/2}
                    (\log\mathcal R)^2
 \tag{A.7}
\]
and a large block for a proportion $(\log\mathcal R)^{-O(1)}$ of these
zeros. The lower bound holds because the coefficients in (A.3) vanish for
$1<n\le U$.
We retain the exponential weight in the coefficients.  The real part of a
zero must also be removed before applying one large-values theorem to the
whole set. Fix the left edge $\sigma$, put $L=\log N$, choose a fixed
$\zeta_0\in C_c^\infty((1/2,2))$ equal to one near
$[1,1+\log(2)/L]$, and set
\[
 \psi_\beta(u)=\zeta_0(u/L)e^{-(\beta-\sigma)u}.
\]
Uniformly for $\sigma\le\beta\le1$,
\[
 \|\psi_\beta^{(j)}\|_1\ll_jL^{1-j},\qquad
 |\widehat\psi_\beta(\xi)|\ll_jL^{1-j}|\xi|^{-j},
 \qquad \|\widehat\psi_\beta\|_1\ll1.                         \tag{A.8}
\]
Fourier inversion gives, up to the fixed Fourier-sign convention,
\[
 \sum_{N<n\le2N}c_ne^{-n/Y}\chi(n)n^{-\beta-i\gamma}
 =\int_{\mathbb R}\widehat\psi_\beta(\xi)
   \sum_{N<n\le2N}c_ne^{-n/Y}\chi(n)
        n^{-\sigma-i(\gamma+2\pi\xi)}\,d\xi.
 \tag{A.9}
\]
Put $H=\mathcal R^\vartheta$. By (A.8), choosing the
integration-by-parts order in terms of $\vartheta$ makes the contribution of
$|\xi|>H$ smaller than the block threshold, uniformly in $\beta$. Indeed,
$|c_n|\le d(n)$ gives the crude uniform estimate
\[
 \sup_t\left|
 \sum_{N<n\le2N}c_ne^{-n/Y}\chi(n)n^{-\sigma-it}
 \right|
 \ll \mathcal R^{o(1)}N^{1-\sigma},
\]
and its product with the Fourier tail from (A.8) is smaller than any prescribed
fixed power of $\mathcal R^{-1}$ after increasing the integration-by-parts
order. Thus every
retained zero chooses some $|\xi_\rho|\le H$ where a common polynomial is
large. Normalize with
\[
 C_N=\max_{N<n\le2N}|c_n|=\mathcal R^{o(1)}
\]
(the block is nonzero), and define the zero-independent polynomial
\[
 D_\chi(t)=\sum_{N<n\le2N}a_n\chi(n)n^{it},
 \qquad a_n=\frac{c_ne^{-n/Y}(N/n)^\sigma}{C_N},
 \qquad |a_n|\le1,
 \tag{A.10}
\]
which satisfies
\[
 |D_\chi(t_j)|\ge N^\sigma\mathcal R^{-o(1)}
 \tag{A.11}
\]
at a shifted ordinate $t_j=-\gamma-2\pi\xi_\rho$. Any unit interval contains
at most $O((H+1)\log\mathcal R)$ of these shifted ordinates, counted with
multiplicity: all their original ordinates lie in an interval of length
$O(H)$, so this follows from (A.5). Greedy selection therefore leaves a
one-separated subset of cardinality
\[
 \gg\#\mathcal Z_I/((H+1)\log\mathcal R).                             \tag{A.12}
\]
It lies in an interval of length $O(T+H)=O(T)$. Translating that interval
multiplies the common coefficients by unimodular factors. This quantifies the
real-part-removal step compressed in \cite[Section~13.1]{GuthMaynard2026}; the fixed factor
$\chi(n)e^{-n/Y}$ remains in the coefficients throughout.

Guth and Maynard's published large-values theorem \cite[Theorem~1.1]{GuthMaynard2026}
applies to every complex coefficient sequence bounded by one.  Hence it
applies to (A.10).  Their power argument in
\cite[Section~13.1]{GuthMaynard2026} has two branches: Theorem~1.1 is used
when the powered length is at most their switching scale, and the usual
discrete mean-value theorem is used in the complementary branch.  Both inputs
allow arbitrary bounded complex coefficients.  In the present twisted setting,
complete multiplicativity gives
\[
 \chi(n_1)\cdots\chi(n_k)=\chi(n_1\cdots n_k),
\]
and, literally,
\[
 D_\chi^k(t)=
 \sum_{N^k<m\le(2N)^k}\chi(m)b_m m^{it},
 \qquad |b_m|\le d_k(m)=\mathcal R^{o(1)}.
\]
Split this support into $O_k(1)$ dyadic blocks. The triangle inequality and
pigeonholing retain one common block for a fixed positive proportion of the
selected ordinates; normalize its divisor-bounded coefficients before
applying \cite[Theorem~1.1]{GuthMaynard2026}. These operations cost only
$\mathcal R^{o(1)}$. By (A.7) and $r\le(\log T)^K$, the polynomial length obeys
$N\le T^{1/2+o_K(1)}$ and $N\ge T^{\kappa+o_K(1)}$. The integer $k$ in
their proof is therefore bounded in terms of $\kappa$; this is enough because
$\kappa$ is fixed before $T\to\infty$. The same power argument gives, for
$7/10\le\sigma\le4/5$,
\[
 \#\mathcal Z_I
 \le T^{\frac{15(1-\sigma)}{3+5\sigma}
          +\vartheta+o_{K,\kappa}(1)}.
 \tag{A.13}
\]
Because
\[
 \frac{15}{3+5\sigma}\le\frac{30}{13}
 \quad\left(\sigma\ge\frac7{10}\right),
 \qquad
 \min_{7/10\le\sigma\le4/5}
 \left(\frac{15}{3+5\sigma}-2\right)(1-\sigma)
 =\frac1{35},
\]
we may first choose $\kappa$ so that $2\kappa<1/70$ and then choose
$\vartheta$ and all remaining subpower losses smaller than the allocated
$\eta/2$. Equations (A.6) and (A.13) prove (A.2) in this range.

For $1/2+\delta\le\sigma\le7/10$, Montgomery's classical family estimate
\cite[Theorem~12.1]{Montgomery1971} gives, for each modulus $r$,
\[
 \sum_{\chi\bmod r}N(\sigma,T,\chi)
 \ll (rT)^{\frac{3(1-\sigma)}{2-\sigma}}
       (\log(rT))^9;
 \tag{A.14}
\]
for the fixed character under consideration, nonnegativity gives explicitly
\[
 N(\sigma,T,\chi)
 \le \sum_{\psi\bmod r}N(\sigma,T,\psi).
\]
Since
\[
 \frac3{2-\sigma}\le\frac{30}{13}
 \quad\left(\sigma\le\frac7{10}\right),
\]
(A.2) follows throughout the asserted strip after taking $T$ sufficiently
large in terms of $K,\delta,\eta$.

Finally $Q^{O(1)}=(\log T)^{O_K(1)}=T^{o_K(1)}$.  Summing (A.2) over the
$O(Q^2)$ original character--modulus pairs and allocating the total subpower
loss to $\eta/2$ proves (A.1).
\end{proof}

\begin{remark}[Scope of the density lemma]
The theorem does not extend the Guth--Maynard estimate to a character family.
It applies that published arbitrary-coefficient theorem separately to each
character and uses that the family is only polylogarithmic. The fixed-character
detector, fourth-moment input, real-part removal, spacing loss, and bounded-power
coefficient normalization are displayed above because they are the uniformity
points on which this reduction depends.
\end{remark}


\begin{thebibliography}{9}
\bibitem{CGL2026}
B. Chen, V. Gupta, and Y. C. Li,
\emph{Large Value Estimates for Dirichlet Polynomials with Characters and Zero Density of Dirichlet $L$-Functions},
arXiv:2507.08296v2 (2026), Theorem 1.2.
\url{https://arxiv.org/pdf/2507.08296v2}.

\bibitem{GafniTao2026}
A. Gafni and T. Tao,
\emph{On the number of exceptional intervals to the prime number theorem in
short intervals},
Ess. Number Th. \textbf{5} (2026), 221--241.
\url{https://arxiv.org/abs/2505.24017v1}.

\bibitem{GuthMaynard2026}
L. Guth and J. Maynard,
\emph{New large value estimates for Dirichlet polynomials},
Ann. of Math. (2) \textbf{203} (2026), no.~2, 623--675.
\url{https://doi.org/10.4007/annals.2026.203.2.6}.

\bibitem{Montgomery1971}
H. L. Montgomery,
\emph{Topics in Multiplicative Number Theory},
Lecture Notes in Mathematics 227, Springer, 1971.

\bibitem{Jutila1977}
M. Jutila,
\emph{On Linnik's constant},
Math. Scand. \textbf{41} (1977), 45--62, Theorem 1, equation (1.7).
\url{https://journals.msp.org/mscand/article/view/1943}.

\bibitem{Khale2024}
T. Khale,
\emph{An explicit Vinogradov--Korobov zero-free region for Dirichlet $L$-functions},
Q. J. Math. \textbf{75} (2024), 299--332, Theorem 1.1.
\url{https://doi.org/10.1093/qmath/haae010}.

\bibitem{Koukoulopoulos}
D. Koukoulopoulos,
\emph{The Distribution of Prime Numbers},
Graduate Studies in Mathematics 203, American Mathematical Society, 2019.
\url{https://dms.umontreal.ca/~koukoulo/documents/publications/primes.pdf}.

\bibitem{MRT2019}
K. Matom\"aki, M. Radziwi\l\l, and T. Tao,
\emph{Correlations of the von Mangoldt and higher divisor functions I: Long shift ranges},
Proc. Lond. Math. Soc. \textbf{118} (2019), 284--350.
\url{https://arxiv.org/abs/1707.01315v3}.

\bibitem{Shiu1980}
P. Shiu,
\emph{A Brun--Titchmarsh theorem for multiplicative functions},
J. Reine Angew. Math. \textbf{313} (1980), 161--170.
\url{https://eudml.org/doc/152201}.
\end{thebibliography}

\end{document}
